Related work is organized by mechanism, not chronology. For each cluster we
state what the cited work \emph{establishes} and what it does \emph{not}, and we
flag the one correction to the project's own notes that surfaced during the
literature scan for this preprint.
\label{sec:rel}

\paragraph{Discrete prompt search (generate-and-search).} \citet{zhou2023ape}
introduced LLM-generated candidate prompts scored on task accuracy
\emph{(automatic prompt engineering)}; \citet{yang2023opro} framed the optimizer
itself as an LLM that proposes candidates from a scored history, and reported the
now-celebrated case where a short phrase (``take a deep breath'')
beat longer instructions. Later scale-limits work \citep{zhang2024reoprolim} shows
the optimizer's size and human-written meta-instructions strongly condition these
results --- a direct warning against letting a scaffold rank without stating its
evaluator. \citet{shin2020autoprompt} searched trigger tokens via gradients;
\citet{fernando2023promptbreeder} evolved prompts; \citet{guo2023evoprompt}
reused evolutionary operators; \citet{pryzant2023protegi} treated natural-language
critiques as ``gradients'' with beam search; \citet{khattab2023dspy} compiled
declarative LM programs; \citet{yuksekgonul2024textgrad} propagated textual
gradients through computation graphs; and \citet{latentprompt2025} searched a
continuous latent space and decoded candidates. What this cluster establishes:
discrete or latent prompt candidates can be searched and scored programmatically.
What it does \emph{not}: it does not establish \emph{which wording does the most
behavioral work per token}, and much of it scores on the same model that proposes,
leaving an evaluator-leakage risk that \citet{zhang2024reoprolim} makes concrete.

\paragraph{Compression and information-theoretic bounds.}
\citet{jiang2023llmlingua} estimate token importance to compress existing prompts
for inference cost; the strongest precedent for this project's distortion-rate
posture is \citet{nagle2024promptrates} (``Fundamental Limits of Prompt
Compression''), which derives the distortion-rate function for black-box prompt
compression as a linear program. That paper establishes a \emph{theoretical}
bound; this toolkit's \texttt{rate\_distortion\_frontier} is an \emph{empirical}
lower envelope, a strictly weaker object. \citet{tishby2000ib} provides the
Information Bottleneck complement (compress input while preserving relevant
information), and \citet{xie2021icl} supplies the theoretical intuition for why
specific wording narrows the effective output distribution. None of these
operations turns a synthetic simulator into evidence about a real model.

\paragraph{Model selection and uncertainty.}
\citet{akinwande2023pacbayes} show that a language-model prior plus a discrete
prompt yields PAC-Bayes generalization bounds usable for model selection --
relevant because selection-confidence needs a penalty for optimizing on the
selection set only. \citet{sabbatella2023bayesprompt} formulate discrete prompt
tuning as combinatorial Bayesian optimization over a continuous surrogate;
\citet{zhou2026bayespo} treat prompts as posterior samples with a language-model
prior (a 2026 result this project cites as \emph{research option}, not evidence).
\citet{angelopoulos2021conformal} give the split-conformal toolkit that the
project's \texttt{conformal\_radius} mirrors in code. These works establish
techniques for honest uncertainty; this project ships the \emph{machinery}, not
the calibrated real-world numbers.

\paragraph{Mechanism: why short words matter.}
\citet{kornblith2019similarity} show that representation-similarity measures such
as CKA do not track task-relevant behavioral similarity --- the paper this
project's ``similarity $\neq$ behavior'' stance leans on. Steering work
\citep{turner2023activation,li2023iti,zou2023repe,panickssery2023caa,meng2022rome}
operates by causal intervention on activations, \emph{not} prompt wording; this
project borrows only the measurement instinct (project a candidate's effect onto
a declared direction) and does not claim access to any production model's
activations. \citet{bartoszcze2025repe} provide the recent survey.

\paragraph{Evaluation robustness.}
\citet{wallace2024instructionhierarchy} show that the priority among
system/developer/user text is itself learned behavior; \citet{li2023promptinj}
show instruction-following degrades under conflicting embedded text (so
``off-target'' should include unauthorized or unsafe behavior, not only semantic
competitors); \citet{wu2023counterfactual} show that apparent reasoning can rely
on default-condition shortcuts, motivating held-out and counterfactual
evaluation rather than more examples from the same template.

\paragraph{Novelty claim (scoped and checkable).}
To our knowledge, no other shipped, reproducible scaffold combines (a) scoring
discrete prompt candidates against \emph{named behavioral attractors}, (b)
empirical distortion-rate \emph{plus} Pareto selection, (c) \emph{mandatory} hard
gates that veto candidates regardless of score, and (d) cooperative-game
attribution (Shapley, interaction, deletion, directional) under one
evidence-gated pipeline whose explicit rule is that embedding-based orderings
must be reconfirmed on real held-out evaluation before use. We assert this
narrowly for the specific combination as implemented in this repository -- not
for the research space generally, and not as any claim of intrinsic novelty for
any single component.

\begin{quote}
\textbf{Correction of record.} The project's own research notes
(\texttt{research-landscape.md} under \texttt{docs/current/}) cite APE as
arXiv:2210.03493. A
live arXiv API check for this preprint shows \texttt{2210.03493} is instead
\textit{Automatic Chain of Thought Prompting in Large Language Models}
\citep{zhuosheng2022autocot}; the correct APE identifier is arXiv:2211.01910
\citep{zhou2023ape}. The present paper cites the corrected ID; the note file
should be updated.
\end{quote}